\hypertarget{sec:hipotesis}{}%
\begin{center}
    {\Large\bfseries 5. FORMULACIÓN DE HIPÓTESIS}
\end{center}
\vspace{1em}

\noindent\textbf{H1a (hipótesis de trabajo --- caracterización funcional):} Los atacantes evaluados mediante el modelo de métricas avanzadas ajustadas por contexto táctico presentan una mayor precisión en la diferenciación de su perfil funcional que cuando son evaluados mediante el modelo tradicional basado en estadísticas basicas (goles y asistencias), en futbolistas con registros estadísticos equivalentes.\par\vspace{0.8em}

\noindent\textbf{H0a (hipótesis nula --- caracterización funcional):} No existe diferencia en la precisión de la diferenciación del perfil funcional de los atacantes entre la evaluación realizada con el modelo de métricas avanzadas ajustadas por contexto táctico y la realizada con el modelo tradicional basado en estadísticas basicas.\par\vspace{1.2em}

\noindent\textbf{H1b (hipótesis de trabajo --- incertidumbre en la evaluación):} La aplicación del modelo de métricas avanzadas contextualizadas presenta un menor grado de incertidumbre y discrepancia en la clasificación del rendimiento real de los atacantes que la evaluación basada exclusivamente en estadísticas tradicionales basicas.\par\vspace{0.8em}

\noindent\textbf{H0b (hipótesis nula --- incertidumbre en la evaluación):} No existe diferencia en el grado de incertidumbre ni en la clasificación del rendimiento real de los atacantes entre la aplicación del modelo de métricas avanzadas contextualizadas y la evaluación basada exclusivamente en estadísticas tradicionales basicas.